\begin{frame}{La brecha tributaria combina incumplimiento y beneficios legales}
\begin{itemize}
    \item Brecha tributaria = brecha de cumplimiento + brecha de política.
    \item La estimación de evasión disponible equivale a 8,6\% del PIB.
    \item El gasto tributario llegó a 8,9\% del PIB en el año gravable 2024.
    \item Bajo el marco tributario de referencia, el recaudo potencial supera ampliamente el observado, aunque no toda la brecha es eliminable.
\end{itemize}
\note{[Sources] MFMP 2026, apéndice, tabla A.2.1. Evasión: Lora et al. (2021), como en la versión fuente.}
\end{frame}

\begin{frame}{El IVA concentra la mayor brecha de política}
\centering
\small
\begin{tabular}{lrrr}
\toprule
\textbf{Concepto (\% del PIB)} & \textbf{IVA} & \textbf{Renta jurídicas} & \textbf{Renta naturales} \\
\midrule
Evasión (2019)                    & 1,8 & 5,7 & 1,1 \\
Gasto tributario (año gravable 2024) & 5,8 & 1,5 & 1,7 \\
Recaudo (2025)                    & 5,3 & 6,0 & 1,5 \\
\bottomrule
\end{tabular}
\vspace{0.25cm}
\scriptsize La evasión y el recaudo provienen de las fuentes de la versión original; el gasto tributario se actualiza con el MFMP 2026.
\note{[Sources] MFMP 2026, apéndice, tabla A.2.1. Evasión: Lora et al. (2021). Recaudo: Balance del GNC, como en la versión fuente.}
\end{frame}

\begin{frame}{Cerrar las brechas exige instrumentos distintos}
\begin{itemize}
    \item La brecha de cumplimiento es principalmente un problema de administración tributaria y control.
    \item La brecha de política es un problema legal y distributivo: reducir beneficios cambia precios relativos y cargas entre hogares y empresas.
    \item El marchitamiento de beneficios puede simplificar el sistema y facilitar una reducción de tarifas, pero enfrenta restricciones jurídicas y políticas.
\end{itemize}
\end{frame}